\documentclass[12pt,a4paper]{article}

% ─────────────────────────────────────────────────────────────────────────────
%  PACKAGES
% ─────────────────────────────────────────────────────────────────────────────
\usepackage[utf8]{inputenc}          % UTF-8 input encoding (fixes Unicode errors)
\usepackage[T1]{fontenc}
\usepackage[a4paper, top=2.2cm, bottom=2.2cm, left=2.2cm, right=2.2cm]{geometry}
\usepackage{xcolor}
\usepackage{fontawesome5}
\usepackage{titlesec}
\usepackage{tcolorbox}
\usepackage{fancyhdr}
\usepackage{lmodern}
\usepackage{microtype}
\usepackage{array}
\usepackage{booktabs}
\usepackage{colortbl}
\usepackage{multirow}
\usepackage{listings}
\usepackage{enumitem}
\usepackage{graphicx}
\usepackage{tikz}
\usepackage{hyperref}
\usepackage{setspace}
\usepackage{amsmath}

\tcbuselibrary{skins, breakable, theorems}
\usetikzlibrary{positioning, shadows}

% Fix headheight warning
\setlength{\headheight}{24.5pt}
\addtolength{\topmargin}{-12.5pt}

% ─────────────────────────────────────────────────────────────────────────────
%  COLOR PALETTE
% ─────────────────────────────────────────────────────────────────────────────
\definecolor{primary}{HTML}{0F4C81}
\definecolor{secondary}{HTML}{2E86C1}
\definecolor{accent}{HTML}{1A5276}
\definecolor{highlight}{HTML}{F39C12}
\definecolor{success}{HTML}{1E8449}
\definecolor{danger}{HTML}{C0392B}
\definecolor{lightbg}{HTML}{EBF5FB}
\definecolor{warningbg}{HTML}{FEF9E7}
\definecolor{successbg}{HTML}{EAFAF1}
\definecolor{codebg}{HTML}{F2F3F4}
\definecolor{linegray}{HTML}{BFC9CA}

% ─────────────────────────────────────────────────────────────────────────────
%  HYPERREF
% ─────────────────────────────────────────────────────────────────────────────
\hypersetup{
    colorlinks=true,
    linkcolor=primary,
    urlcolor=secondary,
    pdfauthor={Dev Jindal et al.},
    pdftitle={Intelligent Credit Risk Prediction System},
}

% ─────────────────────────────────────────────────────────────────────────────
%  PAGE STYLE
% ─────────────────────────────────────────────────────────────────────────────
\pagestyle{fancy}
\fancyhf{}
\fancyhead[L]{\small\color{linegray}\textit{Credit Risk Prediction System}}
\fancyhead[R]{\small\color{linegray}\textit{ML Project Report}}
\fancyfoot[C]{\small\color{linegray}\thepage}
\renewcommand{\headrulewidth}{0.4pt}
\renewcommand{\headrule}{%
  \hbox to\headwidth{\color{linegray}\leaders\hrule height \headrulewidth\hfill}}

% ─────────────────────────────────────────────────────────────────────────────
%  SECTION STYLES
%  NOTE: #1 must NOT appear inside \colorbox / \parbox in \titleformat.
%        The number goes in a pill; the title text flows naturally after it.
% ─────────────────────────────────────────────────────────────────────────────
\titleformat{\section}
  {\normalfont\Large\bfseries\color{primary}}
  {\colorbox{primary}{\color{white}\enspace\thesection\enspace}}
  {0.8em}
  {}
  [{\color{secondary!50}\vspace{3pt}\hrule height 1.5pt}\vspace{2pt}]

\titleformat{\subsection}
  {\normalfont\large\bfseries\color{secondary}}
  {\thesubsection}{0.8em}{}
  [{\color{linegray}\vspace{1pt}\hrule height 0.5pt}\vspace{2pt}]

\titleformat{\subsubsection}
  {\normalfont\normalsize\bfseries\color{accent}}
  {\thesubsubsection}{0.8em}{}

\titlespacing*{\section}{0pt}{2em}{0.6em}
\titlespacing*{\subsection}{0pt}{1.4em}{0.3em}
\titlespacing*{\subsubsection}{0pt}{1em}{0.2em}

% ─────────────────────────────────────────────────────────────────────────────
%  SPACING & LISTS
% ─────────────────────────────────────────────────────────────────────────────
\onehalfspacing
\setlength{\parskip}{0.55em}
\setlength{\parindent}{0pt}
\setlist[itemize]{leftmargin=1.6em, itemsep=0.25em, topsep=0.25em, parsep=0pt}
\setlist[enumerate]{leftmargin=1.8em, itemsep=0.25em, topsep=0.25em, parsep=0pt}

% ─────────────────────────────────────────────────────────────────────────────
%  CUSTOM TCOLORBOX ENVIRONMENTS
% ─────────────────────────────────────────────────────────────────────────────

% Key Insight (blue)
\newtcolorbox{keybox}[1][]{
  enhanced, breakable,
  colback=lightbg, colframe=primary,
  coltitle=white, fonttitle=\bfseries\small,
  title={\faLightbulb\quad Key Insight},
  attach boxed title to top left={yshift=-2mm, xshift=4mm},
  boxed title style={colback=primary, colframe=primary, sharp corners},
  sharp corners=south,
  left=6pt, right=6pt, top=10pt, bottom=6pt,
  #1
}

% Important Note (amber)
\newtcolorbox{warnbox}[1][]{
  enhanced, breakable,
  colback=warningbg, colframe=highlight,
  coltitle=white, fonttitle=\bfseries\small,
  title={\faExclamationTriangle\quad Important Note},
  attach boxed title to top left={yshift=-2mm, xshift=4mm},
  boxed title style={colback=highlight, colframe=highlight, sharp corners},
  sharp corners=south,
  left=6pt, right=6pt, top=10pt, bottom=6pt,
  #1
}

% Strength (green)
\newtcolorbox{successbox}[1][]{
  enhanced, breakable,
  colback=successbg, colframe=success,
  coltitle=white, fonttitle=\bfseries\small,
  title={\faCheckCircle\quad Strength},
  attach boxed title to top left={yshift=-2mm, xshift=4mm},
  boxed title style={colback=success, colframe=success, sharp corners},
  sharp corners=south,
  left=6pt, right=6pt, top=10pt, bottom=6pt,
  #1
}

% ─────────────────────────────────────────────────────────────────────────────
%  TABLE / CODE STYLE
% ─────────────────────────────────────────────────────────────────────────────
\renewcommand{\arraystretch}{1.45}

\lstset{
    backgroundcolor=\color{codebg},
    basicstyle=\ttfamily\small,
    frame=single, rulecolor=\color{linegray},
    breaklines=true,
    numbers=left, numberstyle=\tiny\color{linegray},
    keywordstyle=\color{primary}\bfseries,
    commentstyle=\color{success}\itshape,
    stringstyle=\color{danger},
    showstringspaces=false,
}

% ═══════════════════════════════════════════════════════════════════════════
%  DOCUMENT
% ═══════════════════════════════════════════════════════════════════════════
\begin{document}
\thispagestyle{empty}

% ── COVER PAGE ──────────────────────────────────────────────────────────────
\begin{tikzpicture}[remember picture, overlay]
  \fill[primary] (current page.north west)
    rectangle ([yshift=-5.5cm]current page.north east);
  \fill[secondary!70] (current page.south west)
    rectangle ([yshift=1.2cm]current page.south east);
  \fill[white, opacity=0.07]
    ([xshift=-2.5cm, yshift=-2.5cm]current page.north east) circle (4cm);
  \fill[white, opacity=0.05]
    ([xshift=-1cm, yshift=-1cm]current page.north east) circle (2.5cm);
  \fill[white, opacity=0.07]
    ([xshift=1.5cm, yshift=0.6cm]current page.south west) circle (2.2cm);
\end{tikzpicture}

\vspace*{1.8cm}

\begin{center}
  {\color{white}\fontsize{11}{13}\selectfont
    \textls[200]{MACHINE LEARNING PROJECT REPORT}}\\[0.6cm]
  {\color{white}\fontsize{26}{32}\selectfont\bfseries
    Intelligent Credit Risk}\\[0.2cm]
  {\color{white}\fontsize{26}{32}\selectfont\bfseries
    Prediction System}\\[0.5cm]
  {\color{white!70}\fontsize{12}{15}\selectfont
    Multiclass Borrower Risk Classification using Ensemble ML}
\end{center}

\vspace{3.5cm}

% Info card
\begin{center}
\begin{tcolorbox}[
  enhanced, width=0.84\textwidth,
  colback=white, colframe=linegray,
  arc=6pt, boxrule=0.5pt,
  drop shadow={opacity=0.08},
  left=10pt, right=10pt, top=10pt, bottom=10pt,
]
  \begin{center}
    {\Large\bfseries\color{primary} Project Overview}\\[0.5em]
    {\color{linegray}\rule{0.6\linewidth}{0.4pt}}\\[0.8em]
  \end{center}
  \renewcommand{\arraystretch}{1.5}
  \begin{tabular}{@{} p{3.2cm} p{9cm} @{}}
    \faDatabase\quad\textbf{Datasets}
      & \texttt{case\_study1.csv} (Bank of Baroda) \newline
        \texttt{case\_study2.csv} (CIBIL Bureau) \\
    \faBullseye\quad\textbf{Target}
      & \texttt{Approved\_Flag} --- Multiclass (P1 / P2 / P3 / P4) \\
    \faCogs\quad\textbf{Models}
      & Logistic Regression, Decision Tree, Random Forest, XGBoost \\
    \faWrench\quad\textbf{Tuning}
      & GridSearchCV with 3-fold stratified cross-validation \\
    \faRocket\quad\textbf{Deployment}
      & Gradio web interface for live inference \\
    \faCode\quad\textbf{Stack}
      & Python, scikit-learn, XGBoost, Gradio, pandas \\
  \end{tabular}
\end{tcolorbox}
\end{center}

\vspace{1.8cm}

% Team card
\begin{center}
\begin{tcolorbox}[
  enhanced, width=0.84\textwidth,
  colback=primary!6, colframe=primary,
  arc=6pt, boxrule=0.6pt,
  title={\faUsers\quad Team Members},
  fonttitle=\bfseries\color{white},
  attach boxed title to top left={yshift=-2mm, xshift=4mm},
  boxed title style={colback=primary},
  left=6pt, right=6pt, top=10pt, bottom=6pt,
]
\begin{center}
\renewcommand{\arraystretch}{1.4}
\begin{tabular}{@{} l l @{}}
  \textbf{Dev Jindal} (Team Lead) & Data preprocessing, model development \& coordination \\
  \textbf{Gargi Srivastava}       & Report writing and documentation \\
  \textbf{Agrima Gusain}          & UI design and report writing \\
  \textbf{Bhawana Yadav}          & UI design and implementation \\
  \textbf{Mishti Sharma}          & Demo video and presentation \\
\end{tabular}
\end{center}
\end{tcolorbox}
\end{center}

\vfill
\begin{center}{\color{linegray}\small February 2025}\end{center}

\newpage
{\setlength{\parskip}{0pt}\tableofcontents}
\newpage

% ═══════════════════════════════════════════════════════════════════════════
\section{Abstract}
% ═══════════════════════════════════════════════════════════════════════════

Credit risk assessment is one of the most critical operations for banks and
financial institutions. Before approving a loan, lenders must evaluate whether
a borrower is likely to repay. Incorrect decisions lead to financial losses and
rising default rates --- making an automated, data-driven approach essential.

\begin{keybox}
  This project builds a machine learning credit risk prediction system that
  classifies borrowers into \textbf{four risk tiers} --- P1 (most creditworthy)
  through P4 (highest risk) --- using a merged dataset of internal bank records
  and CIBIL bureau data. The production model is \textbf{Logistic Regression},
  selected for its interpretability and consistent multiclass performance.
\end{keybox}

During initial analysis, the value \texttt{-99999} was discovered in several
numerical columns. Rather than treating these as outliers, investigation confirmed
they were \textbf{sentinel values representing missing data}. Proper handling was
critical to ensure accurate model training.

After cleaning, merging, and performing statistical feature selection
(Chi-square, ANOVA, VIF), four models were trained --- Logistic Regression,
Decision Tree, Random Forest, and XGBoost --- each tuned via GridSearchCV.
A Gradio web interface was developed for real-time borrower risk inference.

\newpage

% ═══════════════════════════════════════════════════════════════════════════
\section{Introduction}
% ═══════════════════════════════════════════════════════════════════════════

\textbf{Credit risk} is the possibility that a borrower will fail to meet
financial obligations as per agreed terms. For banks, credit risk management
is foundational --- lending forms the core of their business, and poor
assessment leads to high defaults, financial instability, and significant losses.

\subsection{Limitations of Traditional Approaches}

Traditional loan approval relied on manual review or fixed rule-based frameworks
--- minimum income thresholds, employment status, basic credit scores.
These systems:

\begin{itemize}
  \item[\textcolor{danger}{\faTimesCircle}]
    Fail to capture \textbf{complex patterns and feature interactions}
  \item[\textcolor{danger}{\faTimesCircle}]
    Are \textbf{time-consuming} and operationally expensive
  \item[\textcolor{danger}{\faTimesCircle}]
    May introduce \textbf{human bias} into decisions
  \item[\textcolor{danger}{\faTimesCircle}]
    Cannot easily scale to thousands of applications daily
\end{itemize}

\subsection{Machine Learning Advantage}

ML algorithms analyze large volumes of historical borrower data and learn hidden
relationships between features --- income, trade line history, loan behavior,
repayment patterns. This project adopts a \textbf{supervised multiclass
classification} approach, predicting one of four risk classes (P1--P4) rather
than a simple binary approve/decline.

\begin{successbox}
  Using stratified splitting, \texttt{class\_weight='balanced'}, and statistical
  feature selection, the system delivers \textbf{robust predictions even for
  minority classes} (P3 and P4) --- the most financially critical cases.
\end{successbox}

\subsection{Key Data Challenge --- Sentinel Values}

\begin{warnbox}
  During EDA, sentinel values of \texttt{-99999.00} were identified in several
  financial columns. These are \textbf{not outliers} --- they are placeholders
  for missing data. Treating them as real values would severely distort the model.
  They were replaced with \texttt{NaN} and handled via standard missing data removal.
\end{warnbox}

\newpage

% ═══════════════════════════════════════════════════════════════════════════
\section{Problem Statement}
% ═══════════════════════════════════════════════════════════════════════════

Financial institutions must evaluate thousands of applications and determine
which borrowers are likely to repay. Manual evaluation is slow and inconsistent.
Rule-based systems fail to capture complex borrower behavior patterns.

\textbf{This project aims to design a machine learning-based credit risk
prediction system} that automatically classifies applicants into one of four
risk tiers with explainable, consistent predictions.

\vspace{0.5em}

\begin{center}
\begin{tcolorbox}[
  enhanced, width=0.92\textwidth,
  colback=primary!6, colframe=primary, arc=6pt,
  title={\faBullseye\quad Four Risk Tiers (Target: \texttt{Approved\_Flag})},
  fonttitle=\bfseries\small\color{white},
  attach boxed title to top left={yshift=-2mm, xshift=4mm},
  boxed title style={colback=primary},
  left=4pt, right=4pt, top=10pt, bottom=4pt,
]
\renewcommand{\arraystretch}{1.5}
\begin{tabular}{>{\bfseries\color{primary}}c p{3.5cm} p{4.5cm} p{3.5cm}}
  \toprule
  \textbf{Class} & \textbf{Risk Profile} & \textbf{Credit Quality} & \textbf{Decision} \\
  \midrule
  P1 & \textcolor{success}{\faCircle}\ Top Tier
     & Excellent history, low risk
     & \textcolor{success}{\textbf{Approve}} \\
  P2 & \textcolor{secondary}{\faCircle}\ Good
     & Minor concerns, manageable
     & \textcolor{secondary}{\textbf{Conditional}} \\
  P3 & \textcolor{highlight}{\faCircle}\ Marginal
     & Notable risk factors present
     & \textcolor{highlight}{\textbf{Manual Review}} \\
  P4 & \textcolor{danger}{\faCircle}\ High Risk
     & High default probability
     & \textcolor{danger}{\textbf{Decline}} \\
  \bottomrule
\end{tabular}
\end{tcolorbox}
\end{center}

\vspace{0.5em}

\textbf{The system focuses on:}

\begin{enumerate}
  \item Merging two data sources (internal bank data + CIBIL bureau data)
  \item Identifying and handling sentinel values (\texttt{-99999})
  \item Statistical feature selection --- Chi-square, VIF, ANOVA
  \item Training and comparing four classification models
  \item Hyperparameter tuning via GridSearchCV with cross-validation
  \item Evaluating with accuracy, macro precision, recall, and F1-score
  \item Deploying a Gradio web interface for live predictions
\end{enumerate}

\newpage

% ═══════════════════════════════════════════════════════════════════════════
\section{Dataset Description}
% ═══════════════════════════════════════════════════════════════════════════

\subsection{Data Sources}

Two datasets are merged on the shared key \texttt{PROSPECTID} via an inner join:

\vspace{0.4em}
\begin{center}
\begin{tabular}{>{\bfseries}p{4cm} p{4.5cm} p{6cm}}
  \toprule
  \textbf{File} & \textbf{Source} & \textbf{Content} \\
  \midrule
  \texttt{case\_study1.csv} & Bank of Baroda (internal) & Demographic \& financial attributes \\
  \texttt{case\_study2.csv} & CIBIL Bureau               & Trade line \& repayment history \\
  \bottomrule
\end{tabular}
\end{center}

\subsection{Feature Overview}

\begin{tcolorbox}[
  enhanced, colback=codebg, colframe=linegray, arc=4pt,
  title={\faTable\quad Key Features in the Merged Dataset},
  fonttitle=\bfseries\small\color{accent},
  attach boxed title to top left={yshift=-2mm, xshift=4mm},
  boxed title style={colback=linegray, colframe=linegray},
  left=4pt, right=4pt,
]
\renewcommand{\arraystretch}{1.3}
\begin{tabular}{>{\bfseries\ttfamily}p{5.6cm} p{2.4cm} p{6.2cm}}
  \textbf{Feature} & \textbf{Type} & \textbf{Description} \\
  \midrule
  NETMONTHLYINCOME           & Numerical   & Net monthly income of borrower \\
  EDUCATION                  & Categorical & Ordinal: SSC to Professional (7 levels) \\
  MARITALSTATUS              & Categorical & Married / Single / Others \\
  GENDER                     & Categorical & M / F \\
  Age\_Oldest\_TL            & Numerical   & Age of oldest trade line (months) \\
  Age\_Newest\_TL            & Numerical   & Age of newest trade line (months) \\
  time\_since\_recent\_payment & Numerical & Months since last payment \\
  time\_since\_recent\_enq   & Numerical   & Months since last credit enquiry \\
  Time\_With\_Curr\_Empr     & Numerical   & Months with current employer \\
  last\_prod\_enq2           & Categorical & Last product enquiry type \\
  first\_prod\_enq2          & Categorical & First product enquiry type \\
  max\_recent\_level\_of\_deliq & Numerical & Max recent delinquency level \\
  recent\_level\_of\_deliq   & Numerical   & Most recent delinquency level \\
\end{tabular}
\end{tcolorbox}

\subsection{Data Quality Issues}

\begin{warnbox}
  \textbf{Sentinel values (\texttt{-99999})} were found across multiple numerical
  columns. These represent \emph{missing data}, not extreme financial measurements.
  Mishandling them as real values would introduce severe distortions into the model.

  \vspace{0.4em}
  \textbf{Resolution:} All sentinel values were replaced with \texttt{NaN} via
  \texttt{clean\_df1(sentinel=-99999.0)}, followed by row-level removal.
  Only $\approx$\,40 rows were dropped --- a negligible fraction of the total dataset.
  Two CIBIL columns with excessive missingness were dropped entirely.
\end{warnbox}

\newpage

% ═══════════════════════════════════════════════════════════════════════════
\section{Data Preprocessing}
% ═══════════════════════════════════════════════════════════════════════════

\subsection{Step-by-Step Pipeline}

\begin{enumerate}

  \item \textbf{Sentinel Replacement}\\
  All \texttt{-99999} values replaced with \texttt{NaN} using
  \texttt{clean\_df1(df, sentinel=-99999.0)}. Standard missing-value
  handling is then applied uniformly downstream.

  \item \textbf{Row Removal}\\
  Rows with any remaining \texttt{NaN} values removed. Impact: $\approx$40
  rows dropped --- negligible relative to dataset size.

  \item \textbf{Column Dropping}\\
  Two CIBIL columns with excessive missingness dropped before merging --- too
  sparse to impute reliably.

  \item \textbf{Inner Join}\\
  \texttt{df1.merge(df2, on='PROSPECTID', how='inner')} --- only records present
  in both data sources are retained, ensuring data completeness.

  \item \textbf{Identifier Removal}\\
  \texttt{PROSPECTID} dropped after joining (carries no predictive value).

  \item \textbf{Feature Encoding}

  \begin{tcolorbox}[
    colback=lightbg, colframe=secondary, arc=4pt,
    left=6pt, right=6pt, top=4pt, bottom=4pt,
  ]
  \renewcommand{\arraystretch}{1.4}
  \begin{tabular}{@{} p{2.8cm} p{10.5cm} @{}}
    \textbf{Ordinal}
      & \texttt{EDUCATION} mapped to integers 1--7
        (SSC=1, 12TH=2, Graduate=3, Under Graduate=4,
         Post-Graduate=5, Professional=6, Others=7) \\
    \textbf{One-hot}
      & \texttt{MARITALSTATUS}, \texttt{GENDER},
        \texttt{last\_prod\_enq2}, \texttt{first\_prod\_enq2}
        encoded via \texttt{pd.get\_dummies()} \\
    \textbf{Label Enc.}
      & Target \texttt{Approved\_Flag} encoded to integers via
        \texttt{LabelEncoder}; inverse-transformed at inference time \\
  \end{tabular}
  \end{tcolorbox}

  \item \textbf{Statistical Feature Selection}

  Three methods applied in sequence:
  \begin{itemize}
    \item \textbf{Chi-square test} --- categorical features: retain at $p < 0.05$
    \item \textbf{VIF} (Variance Inflation Factor) --- iterative removal of
          multicollinear numerical features above threshold 6.0
    \item \textbf{ANOVA F-test} --- numerical features: retain at $p < 0.05$
  \end{itemize}

  \item \textbf{Scaling and Train/Test Split}

  \begin{successbox}
    \textbf{Zero data leakage:} \texttt{StandardScaler} was \emph{fit on the
    training set only} and \emph{transform-only} applied to the test set.
    An 80/20 split with \texttt{stratify=y} preserved the P1--P4 class
    distribution in both sets. \texttt{RANDOM\_STATE=42} enforces reproducibility.
  \end{successbox}

\end{enumerate}

\newpage

% ═══════════════════════════════════════════════════════════════════════════
\section{Modeling Methodology}
% ═══════════════════════════════════════════════════════════════════════════

Since \texttt{Approved\_Flag} has four classes, this is a \textbf{multiclass
classification} problem. Four models were trained and compared via a shared
\texttt{train\_with\_grid\_search()} utility (\texttt{n\_jobs=-1}, \texttt{verbose=1}).

\subsection{Logistic Regression --- Production Model}

\begin{tcolorbox}[
  enhanced, colback=successbg, colframe=success, arc=5pt,
  title={\faRocket\quad Selected as Production Model},
  fonttitle=\bfseries\small\color{white},
  attach boxed title to top left={yshift=-2mm, xshift=4mm},
  boxed title style={colback=success},
  left=6pt, right=6pt, top=10pt, bottom=6pt,
]
  Logistic Regression was selected over higher-accuracy ensemble models because:
  \begin{itemize}
    \item[\faCheck] Produces \textbf{well-calibrated probabilities}
          --- essential for risk-tier thresholds
    \item[\faCheck] \textbf{Interpretable} --- auditable in a
          lending / regulatory context
    \item[\faCheck] \textbf{Consistent macro-level performance}
          across all four risk classes
  \end{itemize}
\end{tcolorbox}

\textbf{Hyperparameter search space:}
\begin{itemize}
  \item \texttt{C}: [0.01, 0.1, 1, 10] \quad
        \texttt{solver}: [lbfgs, saga] \quad
        \texttt{max\_iter}: [1000]
  \item \texttt{class\_weight}: \texttt{'balanced'} (fixed) \quad
        \texttt{multi\_class}: \texttt{'multinomial'} (fixed)
\end{itemize}

\subsection{Decision Tree}

A Decision Tree captures nonlinear feature interactions.
\texttt{class\_weight='balanced'} up-weights minority classes.

\textbf{Grid:} \texttt{max\_depth}: [3, 5, 7] \quad
\texttt{min\_samples\_split}: [2, 5, 10] \quad
\texttt{criterion}: [gini, entropy]

\subsection{Random Forest}

An ensemble of 100--200 trees reduces overfitting vs.\ a single Decision Tree.

\textbf{Grid:} \texttt{n\_estimators}: [100, 200] \quad
\texttt{max\_depth}: [5, 10, None] \quad
\texttt{min\_samples\_split}: [2, 5]

\subsection{XGBoost}

Gradient-boosted trees with \texttt{objective='multi:softprob'} and
\texttt{eval\_metric='mlogloss'} for proper multiclass probability output.

\textbf{Grid:} \texttt{n\_estimators}: [100, 200] \quad
\texttt{max\_depth}: [3, 5] \quad
\texttt{learning\_rate}: [0.05, 0.1] \quad
\texttt{subsample}: [0.8, 1.0]

\subsection{GridSearchCV Configuration}

\begin{center}
\begin{tabular}{>{\bfseries}p{3.5cm} p{9cm}}
  \toprule
  \textbf{Setting} & \textbf{Value} \\
  \midrule
  Strategy     & Exhaustive grid search across all combinations \\
  CV Folds     & 3 (stratified) \\
  Scoring      & Accuracy \\
  Parallelism  & \texttt{n\_jobs=-1} (all available cores) \\
  Best model   & \texttt{best\_estimator\_} stored in \texttt{model\_registry} \\
  \bottomrule
\end{tabular}
\end{center}

\newpage

% ═══════════════════════════════════════════════════════════════════════════
\section{Model Evaluation}
% ═══════════════════════════════════════════════════════════════════════════

All four models were evaluated on the held-out 20\% test set.

\subsection{Metrics}

\begin{center}
\begin{tabular}{>{\bfseries}p{3.5cm} p{10cm}}
  \toprule
  \textbf{Metric} & \textbf{Description} \\
  \midrule
  Accuracy         & Overall proportion of correct predictions \\
  Macro Precision  & Average precision across P1, P2, P3, P4 (equal weight per class) \\
  Macro Recall     & Average recall across all four classes \\
  Macro F1         & Harmonic mean of macro precision and recall \\
  Per-class P/R/F1 & Individual breakdown for each of P1, P2, P3, P4 \\
  \bottomrule
\end{tabular}
\end{center}

\begin{warnbox}
  In credit risk, \textbf{per-class metrics matter more than overall accuracy}.
  Classifying a P4 borrower (high risk) as P1 (top tier) is a catastrophic error
  --- far more damaging than the reverse. Per-class precision and recall expose
  exactly this kind of inter-class failure.
\end{warnbox}

\subsection{Model Comparison}

Results from all four models were aggregated into a comparison table sorted by
accuracy (descending). A full \texttt{classification\_report} was printed per
model with readable P1--P4 labels and per-class support counts.

\subsection{Production Model Selection}

\begin{keybox}
  \textbf{Logistic Regression} was selected as the production model.
  While ensemble methods may achieve marginally higher accuracy,
  Logistic Regression delivers:
  \begin{itemize}
    \item[\faAngleRight] \textbf{Calibrated probabilities}
          --- critical for defining risk-tier cut-offs
    \item[\faAngleRight] \textbf{Interpretable coefficients}
          --- required for regulatory lending compliance
    \item[\faAngleRight] \textbf{Consistent recall on P3/P4}
          --- the minority, high-stakes borrower classes
  \end{itemize}
  This choice is \textbf{intentional and documented} in the \texttt{predict()} function.
\end{keybox}

\newpage

% ═══════════════════════════════════════════════════════════════════════════
\section{Model Optimization}
% ═══════════════════════════════════════════════════════════════════════════

\subsection{Hyperparameter Tuning}

GridSearchCV tested all parameter combinations defined in the central
\texttt{CONFIG} dictionary. Best parameters were automatically selected
via \texttt{best\_estimator\_} and stored in \texttt{model\_registry}.

\subsection{Cross-Validation}

3-fold cross-validation was applied during GridSearchCV, ensuring performance
estimates are not dependent on a single split. Each fold preserved the class
distribution via internal stratification.

\subsection{Class Imbalance Handling}

\begin{successbox}
  \textbf{Three strategies} combined to handle class imbalance across P1--P4:
  \begin{enumerate}
    \item \texttt{class\_weight='balanced'} on LR, DT, RF
          --- up-weights minority classes P3 and P4 during training
    \item \texttt{stratify=y} during train/test split
          --- preserves the class ratio in both sets
    \item \textbf{Macro F1} reported separately
          --- not fooled by majority class dominance in accuracy
  \end{enumerate}
\end{successbox}

\subsection{Feature Refinement}

Statistical feature selection (Chi-square, VIF, ANOVA) reduced noise,
removed multicollinear features, and retained only the most statistically
significant predictors, improving model generalization.

\subsection{Inference Pipeline Consistency}

\begin{keybox}
  The \texttt{predict(raw\_input: dict)} function mirrors the training
  pipeline \textbf{step-for-step}:
  \begin{enumerate}
    \item Build single-row DataFrame from \texttt{raw\_input}
    \item Ordinal-encode \texttt{EDUCATION} via \texttt{CONFIG['EDUCATION\_MAP']}
    \item One-hot encode nominal columns
    \item \texttt{reindex()} to align with training feature set
          (fills unseen OHE columns with 0)
    \item Scale continuous features using the \emph{training-fitted} scaler
    \item Predict with \texttt{lr\_model} and inverse-transform the label
  \end{enumerate}
  This guarantees \textbf{zero training--inference mismatch}.
\end{keybox}

\newpage

% ═══════════════════════════════════════════════════════════════════════════
\section{Gradio Web Interface}
% ═══════════════════════════════════════════════════════════════════════════

A Gradio-based interface provides real-time credit risk inference for
non-technical users, organized into three columns.

\subsection{Input Fields (13 Total)}

\begin{center}
\begin{tcolorbox}[
  enhanced, colback=lightbg, colframe=secondary, arc=5pt,
  title={\faDesktop\quad UI Layout --- Three-Column Design},
  fonttitle=\bfseries\small\color{white},
  attach boxed title to top left={yshift=-2mm, xshift=4mm},
  boxed title style={colback=secondary},
  left=4pt, right=4pt, top=10pt, bottom=4pt,
]
\renewcommand{\arraystretch}{1.35}
\begin{tabular}{>{\bfseries}p{3.2cm} p{4.5cm} p{5.5cm}}
  \toprule
  \textbf{Column} & \textbf{Field} & \textbf{Widget} \\
  \midrule
  \multirow{4}{*}{\color{primary}Personal Info}
    & Education             & Dropdown (7 levels) \\
    & Marital Status        & Dropdown (3 choices) \\
    & Gender                & Dropdown (M / F) \\
    & Net Monthly Income    & Number field (Rs.) \\
  \midrule
  \multirow{4}{*}{\color{secondary}Credit Profile}
    & Last Product Enquiry  & Dropdown (6 types) \\
    & First Product Enquiry & Dropdown (6 types) \\
    & Age of Oldest TL      & Slider (0--300 months) \\
    & Age of Newest TL      & Slider (0--120 months) \\
  \midrule
  \multirow{5}{*}{\color{accent}Delinquency}
    & Months Since Payment  & Slider (0--120) \\
    & Months Since Enquiry  & Slider (0--60) \\
    & Months With Employer  & Slider (0--300) \\
    & Max Delinquency Level & Slider (0--10) \\
    & Recent Deliq. Level   & Slider (0--10) \\
  \bottomrule
\end{tabular}
\end{tcolorbox}
\end{center}

\subsection{Outputs}

\begin{center}
\renewcommand{\arraystretch}{1.4}
\begin{tabular}{>{\bfseries}p{3cm} p{10.5cm}}
  \toprule
  \textbf{Output} & \textbf{Description} \\
  \midrule
  Prediction
    & Risk class with color indicator and lending recommendation:\\
    & \quad \textcolor{success}{\faCircle}\enspace P1 --- Top tier (approve)\\
    & \quad \textcolor{secondary}{\faCircle}\enspace P2 --- Good (approve with conditions)\\
    & \quad \textcolor{highlight}{\faCircle}\enspace P3 --- Marginal (manual review)\\
    & \quad \textcolor{danger}{\faCircle}\enspace P4 --- High risk (decline)\\
  \midrule
  Probabilities
    & Percentage confidence for all four classes (P1--P4) \\
  \bottomrule
\end{tabular}
\end{center}

\begin{successbox}
  The Gradio wrapper calls the \textbf{full} \texttt{predict()} function ---
  all encoding and scaling steps are applied correctly.
  No pipeline step is bypassed.
\end{successbox}

\newpage

% ═══════════════════════════════════════════════════════════════════════════
\section{Conclusion}
% ═══════════════════════════════════════════════════════════════════════════

This project built a machine learning credit risk prediction system from two
merged real-world data sources --- internal bank records and CIBIL bureau data.
The complete workflow covered data loading, EDA, sentinel value handling,
preprocessing, statistical feature selection, multiclass model training,
hyperparameter optimization, evaluation, and Gradio deployment.

\begin{keybox}
  \textbf{Key achievements:}
  \begin{itemize}
    \item[\faCheck] Four models trained and compared on a 4-class risk problem
    \item[\faCheck] Sentinel values correctly identified and handled
    \item[\faCheck] Zero data leakage --- scaler fitted on training set only
    \item[\faCheck] Stratified splitting preserves class balance
    \item[\faCheck] Production \texttt{predict()} function mirrors
          the training pipeline step-for-step
    \item[\faCheck] Live Gradio interface for non-technical users
  \end{itemize}
\end{keybox}

Logistic Regression was selected as the production model for its interpretability,
calibrated probabilities, and consistent performance across all four risk classes.
The use of \texttt{class\_weight='balanced'} and statistical feature selection
contributed to robust generalization.

This system demonstrates that machine learning can \textbf{meaningfully improve
lending decision quality} --- replacing slow, biased rule-based systems with
accurate, nuanced, and explainable risk stratification.

\newpage

% ═══════════════════════════════════════════════════════════════════════════
\section{Future Work}
% ═══════════════════════════════════════════════════════════════════════════

\begin{tcolorbox}[
  enhanced, colback=warningbg, colframe=highlight, arc=5pt,
  title={\faForward\quad Recommended Improvements},
  fonttitle=\bfseries\small\color{white},
  attach boxed title to top left={yshift=-2mm, xshift=4mm},
  boxed title style={colback=highlight},
  left=6pt, right=6pt, top=10pt, bottom=6pt,
]
\begin{enumerate}
  \item \textbf{Model Persistence} --- Save model, scaler, and encoder
        as a \texttt{joblib} bundle; globals are lost on kernel restart

  \item \textbf{GridSearch Scoring} --- Switch from \texttt{accuracy}
        to \texttt{macro-f1} for better optimization on minority classes (P3, P4)

  \item \textbf{Confusion Matrix} --- Add to evaluation; essential for
        understanding inter-class misclassification costs

  \item \textbf{XGBoost Imbalance Handling} --- Add sample weights
        (equivalent of \texttt{class\_weight='balanced'} for XGBoost)

  \item \textbf{CV Folds} --- Raise from 3 to 5 for more stable estimates

  \item \textbf{Explainability (XAI)} --- Integrate SHAP or LIME for
        regulatory-grade model interpretability

  \item \textbf{Package Versioning} --- Pin \texttt{scikit-learn},
        \texttt{xgboost}, \texttt{gradio} in \texttt{requirements.txt}

  \item \textbf{Cloud Deployment} --- Deploy via Hugging Face Spaces
        or Render for persistent public access

  \item \textbf{Additional Models} --- Experiment with LightGBM,
        CatBoost, and stacked ensembles
\end{enumerate}
\end{tcolorbox}

\end{document}